\begin{IEEEbiography}[{\includegraphics[width=1in,height=1.25in,clip,keepaspectratio]{photo/angli.png}}]
{Ang Li} received the B.S. degree in instrumentation science and technology from Jilin University, China, in 2020, where he is currently pursuing the Ph.D. degree. His research interests include physics-informed neural networks for sensor signal processing, nonlinear compensation of geophysical instruments, embedded AI implementation for sensing systems, and electrochemical seismometer technology.
\end{IEEEbiography}


\begin{IEEEbiography}[{\includegraphics[width=1in,height=1.25in,clip,keepaspectratio]{photo/hongyuanyang.jpg}}]
{Hongyuan Yang} received the Ph.D. degree in geophysical exploration and information technology from Jilin University, Changchun, China, in 2009. He was a Visiting Scholar with the University of Houston, Houston, TX, USA, from 2017 to 2018. He is currently a Professor with the College of Instrumentation and Electrical Engineering, Jilin University. His research interests include development of seismic exploration instrument, wireless seismic data acquisition, and seismic data processing.
\end{IEEEbiography}


\begin{IEEEbiography}[{\includegraphics[width=1in,height=1.25in,clip,keepaspectratio]{photo/fanzheng.jpg}}]
{Fan Zheng} received the Ph.D. degree in circuits and systems from Jilin University, Changchun, China, in 2008. He was a Visiting Scholar with Uppsala University, Uppsala, Sweden, from 2016 to 2017. He is currently an Associate Professor with the College of Instrumentation and Electrical Engineering, Jilin University. His research interests include development of seismic exploration instrument, wireless seismic data acquisition, and seismic data processing.
\end{IEEEbiography}


\begin{IEEEbiography}[{\includegraphics[width=1in,height=1.25in,clip,keepaspectratio]{photo/huaizhuzhang.png}}]
{Huaizhu Zhang} received the B.S. degree in Mechatronics Engineering in 1998 from Jilin University of Technology, and the Ph.D. degree in 2008 from Changchun Institute of Optics and Machinery, Chinese Academy of Sciences. He was a Visiting Scholar with the University of Houston, Houston, TX, USA, from 2018 to 2019. He is currently an Associate Professor with the College of Instrumentation and Electrical Engineering, Jilin University. His research interests include seismic instruments, data acquisition and signal processing.
\end{IEEEbiography}


\begin{IEEEbiography}[{\includegraphics[width=1in,height=1.25in,clip,keepaspectratio]{photo/linhangzhang.jpg}}]
{Linhang Zhang} received the Ph.D. degree in Earth exploration \& information technology from Jilin University, Changchun, China, in 2007. He was a Visiting Scholar with the University of Houston, Houston, TX, USA, from 2019 to 2020. He is currently an Associate Professor with the College of Instrumentation and Electrical Engineering, Jilin University. His research interests include development of seismic exploration instrument, wireless seismic data acquisition, and seismic data processing.
\end{IEEEbiography}


\begin{IEEEbiography}[{\includegraphics[width=1in,height=1.25in,clip,keepaspectratio]{photo/xunqiantong.jpg}}]
{Xunqian Tong} received the Ph.D. degree from Jilin University, China, in 2016. He was a cotutelle student in Macquarie University, Australia under the CSC sponsor from 2014 to 2016. He is currently an Associate Professor with the College of Instrumentation \& Electrical Engineering, Jilin University. His research interests include IOT protocol design and LPWAN performance analysis.
\end{IEEEbiography}


\begin{IEEEbiography}[{\includegraphics[width=1in,height=1.25in,clip,keepaspectratio]{photo/canliu.png}}]
{Can Liu} is currently pursuing the Ph.D. degree at Jilin University. He received the Master's degree in Instrument Science and Technology from North University of China in 2022. His research interests include sensor signal processing and nonlinear compensation techniques for geophysical instruments.
\end{IEEEbiography}


\begin{IEEEbiography}[{\includegraphics[width=1in,height=1.25in,clip,keepaspectratio]{photo/ruojinli.jpg}}]
{Ruojin Li} received the B.S. degree in instrumentation science and technology from Jilin University, China, in 2024, where she is currently pursuing the Ph.D. degree. Her research interests include application of sensor signal processing, zero-phase correction of geophysical instruments ,and force feedback technology in seismic sensing systems.
\end{IEEEbiography}


\end{document}


